% Part: lambda-calculus
% Chapter: syntax
% Section: de-bruijn

\documentclass[../../../include/open-logic-section]{subfiles}

\begin{document}

\olfileid{lam}{syn}{deb}

\olsection{د دې بروين شاخص}

$\alpha$-معادلتوب ډېر طبيعي دے، ځکه $\alpha$-معادل ترمونه «يوه معنا»
لري۔ په حقيقت کښې د لامبډا ترمونو داسې نحو هم ورکول کېدے شي چې
هغه ترمونه سره نۀ بېلوي کوم چې د $\alpha$-بدلون په وسيله يو بل ته
اوړي۔ تر ټولو پېژندل شوې لاره د متغيرونو پر ځاے د هغو د
\emph{دې بروين شاخص} کارول دي۔

کله چې $\lambd[x][M]$ ليکو، په څرګند ډول وايو چې $x$ د تابع
پاراميټر دے، څو په $M$ کښې د~$x$ په کارولو همدې پاراميټر ته
اشاره وکړو۔ خو د دې
بروين په شاخصي بڼه کښې پاراميټرونه نومونه نۀ لري؛ د تابع په بدنه
کښې د هغوی يادونه په داسې شمېر کېږي چې د کارېدونکي مورد او د هغه
د تړونکي تجريد تر منځ د تجريدونو شمېر ښيي۔ د بېلګې په توګه،
$\lambd[x][\lambd[y][y x]]$ ته وګورئ: دباندنے تجريد متغير~$x$
تړي؛ دنننے تجريد متغير~$y$ تړي؛ فرعي ترم $y x$ د دننني تجريد
په ساحه کښې دے۔ د $y$ او د هغه د تړونکي $\lambd[y]$ تر منځ
هېڅ تجريد نشته، خو د $x$ او د هغه د تړونکي $\lambd[x]$ تر منځ
يو تجريد دے۔ نو د $y x$ دپاره $0\, 1$ ليکو، او د ټول ترم
دپاره $\lambd[][\lambd[][01]]$ ليکو۔ دلته وروستۍ دوه څنګ په
څنګ نښې تطبيق ښيي، يوه لسګونه شمېره نۀ ښيي۔

\begin{defn}
  د دې بروين ترمونه په استقرايي ډول داسې تعريفېږي:
  \begin{enumerate}
  \item $n$، چې $n$ هر طبيعي عدد وي۔
  \item $PQ$، چې $P$ او $Q$ دواړه د دې بروين ترمونه وي۔
  \item $\lambd[][N]$، چې $N$ د دې بروين ترم وي۔
  \end{enumerate}
\end{defn}

د عادي لامبډا ترمونو د دې بروين په شاخصي ترمونو بدلون په رسمي ډول
داسې دے:
\begin{defn}
  \begin{align*}
    F_\Gamma(x) &= \Gamma(x) \\
    F_\Gamma(PQ) &= F_\Gamma(P)F_\Gamma(Q) \\
    F_\Gamma(\lambd[x][N]) &= \lambd[][F_{x,\Gamma}(N)]
  \end{align*}
  چې $\Gamma$ د متغيرونو له صفره شاخص لرونکے لړ دے، او
  $\Gamma(x)$ په $\Gamma$ کښې د متغير~$x$ ځای ښيي۔ د بېلګې په
  توګه، که $\Gamma$ د $x,y,z$ لړ وي، نو $\Gamma(x)$، $0$
  او $\Gamma(z)$، $2$ دے۔
  \textbf{د ژباړې د نسخې سپيناوی (OLLAM-025):} که يو نوم په
  لړ کښې څو ځله راغلے وي، له سر څخه د هغه لومړے ځای اخلو؛
  همدا د دننني هم‌نومي تړونکي شاخص ټاکي۔ سرچينه د تکراري نومونو
  دا قاعده په څرګنده نۀ وايي۔

  $x,\Gamma$ هغه لړ ښيي چې د $\Gamma$ سر ته $x$ ورزيات کړے
  شي؛ د بېلګې په توګه، د مخکنۍ بېلګې د $\Gamma$ په دوام کښې
  $w,\Gamma$ د $w,x,y,z$ لړ دے۔
\end{defn}

له دې بروين ترم څخه د معياري لامبډا ترم بېرته ترلاسه کول داسې
کېږي:

\begin{defn}
  \begin{align*}
    G_\Gamma(n) &= \Gamma[n] \\
    G_\Gamma(PQ) &= G_\Gamma(P) G_\Gamma(Q) \\
    G_\Gamma(\lambd[][N]) &= \lambd[x][G_{x,\Gamma}(N)]
  \end{align*}
  چې $\Gamma$ بيا د متغيرونو له صفره شاخص لرونکے لړ دے، او
  $\Gamma[n]$ د~$n$ په ځای کښې متغير ښيي۔ د بېلګې په توګه،
  که $\Gamma$ د $x,y,z$ لړ وي، نو $\Gamma[1]$، $y$ دے۔

  په وروستۍ معادله کښې متغير~$x$ داسې هر متغير ټاکل کېږي چې
  په $\Gamma$ کښې نۀ وي۔
  \textbf{د ژباړې د نسخې سپيناوی (OLLAM-026):} دا بېرته
  اخيستنه يوازې هغه وخت ټاکلې ده چې ياد شوے ځای په اړوند
  لړ کښې موجود وي؛ سرچينه دغه د تعريف ساحې شرط په څرګنده
  نۀ بيانوي۔
\end{defn}

دلته څو پايلې بې له ثبوته راوړو:

\begin{prop}
  که $M \aconvone M'$ وي، او $\Gamma$ داسې هر لړ وي چې
  $\FV{M}$ پکې شامل وي، نو $F_\Gamma(M) \eqs F_\Gamma(M')$۔
\end{prop}

\end{document}
